% GENERATED FILE. Edit canonical lessons, then run npm run generate:books.
% canonical-source-hash: fnv1a64:ffa73453acdb12d1
% canonical-lessons: GU-C282-sauthi, GU-W52-au-sign, GU-C282-kartan, GU-C282-sudhi, GU-C282-be-varasthi, GU-C282-rtu, GU-W53-vocalic-r-letter, GU-R282-comparing-and-time

\chapter{Comparing, and how long}
\label{ch:gu-comparing-and-how-long}
\hlchaptermodality{282}

\begin{chapteropening}
\textbf{By the end of this chapter:} \emph{I can compare two things (X Y kartā̃ moṭũ che), say what is most (sauthī), say how long and since when (be varasthī, savārthī sā̃j sudhī), name the seasons with ṛtu, and read and write \gu{ૌ} and \gu{ઋ}.}

Complete GU-R282-comparing-and-time: compare things, and say how long and since when.
\end{chapteropening}

\section[\texorpdfstring{\gu{સૌથી} (sauthī)}{સૌથી (sauthī)}]{\gu{સૌથી} (sauthī) — most; of all}

\label{lesson:GU-C282-sauthi}



\begin{quote}
Before the new one: say \textquotedblleft{}I want to become a doctor\textquotedblright{}. (\emph{māre dāktar banvũ che}.) Then \textquotedblleft{}small\textquotedblright{}, \textquotedblleft{}more\textquotedblright{} and \textquotedblleft{}to like\textquotedblright{}. (\textbf{\gu{નાનું}}, \textbf{\gu{વધારે}}, \textbf{\gu{ગમવું}}.)
\end{quote}

\subsection*{You'll want to know: \gu{સૌથી}}
\textbf{\gu{સૌથી}} — \emph{sauthī} — \textquotedblleft{}most; of all\textquotedblright{}. Put it before a describing word and it becomes the most of all:

\begin{quote}
\textbf{\gu{મીરા} \gu{સૌથી} \gu{નાની} \gu{છે}.}
\end{quote}

\begin{quote}
\emph{Mīrā sauthī nānī che.} — \textquotedblleft{}Mira is the youngest.\textquotedblright{}
\end{quote}

\textbf{\gu{નાની}} follows Mira; \textbf{\gu{સૌથી}} stays as it is. With \textbf{\gu{વધારે}} it is \textquotedblleft{}most of all\textquotedblright{}:

\begin{quote}
\textbf{\gu{મને} \gu{ચોમાસું} \gu{સૌથી} \gu{વધારે} \gu{ગમે} \gu{છે}.}
\end{quote}

\begin{quote}
\emph{mane comāsũ sauthī vadhāre game che.} — \textquotedblleft{}I like the monsoon most of all.\textquotedblright{}
\end{quote}

\begin{cousinweb}
\textbf{\gu{સૌથી}} (\emph{sauthī}) is built from \textbf{\gu{સૌ}}, \textquotedblleft{}all, everyone\textquotedblright{}, and \textbf{-\gu{થી}}, \textquotedblleft{}from\textquotedblright{}: \textquotedblleft{}from all\textquotedblright{}, so \textquotedblleft{}more than all\textquotedblright{}. CDIAL lists Gujarati \emph{sau}, \textquotedblleft{}wholly\textquotedblright{}, under Sanskrit \emph{sarva}, \textquotedblleft{}all, whole\textquotedblright{} (CDIAL 13276).
\end{cousinweb}

\subsection*{Your turn}
Say these aloud:

\begin{itemize}
\raggedright
  \item \emph{sauthī}, then \emph{sauthī nānũ}
  \item \textquotedblleft{}Mira is the youngest\textquotedblright{} — \emph{Mīrā sauthī nānī che}
  \item \textquotedblleft{}I like the monsoon most of all\textquotedblright{} — \emph{mane comāsũ sauthī vadhāre game che}
\end{itemize}

\begin{tcolorbox}[breakable,colback=teal!4,colframe=teal!35!black,title={Before you move on}]
What does \textbf{\gu{સૌથી}} mean? (\textquotedblleft{}Most; of all\textquotedblright{}.) \textquotedblleft{}The biggest\textquotedblright{}, of a house? (\textbf{\gu{સૌથી} \gu{મોટું}}.)
\end{tcolorbox}
\section[\texorpdfstring{\gu{ૌ} (au sign)}{ૌ (au sign)}]{\gu{ૌ} (au sign) — the au sign in \gu{સૌથી}}

\label{lesson:GU-W52-au-sign}



\begin{quote}
Before the new one: say \textbf{\gu{સૌથી}} (\emph{sauthī}), \textquotedblleft{}most of all\textquotedblright{}. Then write \textbf{\gu{ૈ}}, the \emph{ai} sign, and \textbf{\gu{ો}}, the \emph{o} sign, once each.
\end{quote}

\subsection*{Script you'll notice: \gu{ૌ}}
\textbf{\gu{ૌ}} — the \emph{au} sign: the \textbf{\gu{ા}} bar on the right plus two stacked strokes above. It sits on \textbf{\gu{સ}} in \textbf{\gu{સૌથી}}.

\noindent
\begin{tabularx}{\linewidth}{@{}>{\raggedright\arraybackslash}X>{\raggedright\arraybackslash}X@{}}
\toprule
\textbf{one stroke above} & \textbf{two strokes above} \\
\midrule
\textbf{\gu{ે}} \emph{e} & \textbf{\gu{ૈ}} \emph{ai}, as in \textbf{\gu{પૈસા}} \\
\textbf{\gu{ો}} \emph{o}: \textbf{\gu{ા}} with \textbf{\gu{ે}} above & \textbf{\gu{ૌ}} \emph{au}: \textbf{\gu{ા}} with \textbf{\gu{ૈ}} above \\
\bottomrule
\end{tabularx}

It is written after its consonant: \textbf{\gu{સ}} first, then the bar to its right, then the lower flag, then the upper flag.

\subsection*{Writing — notice and trace}
\hlblockfigure{\detokenize{figures/GU-W52-au-sign-filmstrip.pdf}}{How \gu{ૌ} is written, stroke by stroke}

Follow the numbered strip:

\begin{enumerate}
\raggedright
  \item start at the top of the bar and draw it straight down
  \item without lifting, turn into the short foot to the right
  \item lift, then start at the lower flag's upper-left tip and arc it to the right
  \item without lifting, curve it down to the right
  \item lift, then start at the upper flag's upper-left tip and arc it to the right
  \item without lifting, curve it down to the right — and only now lift
\end{enumerate}

Three strokes — six movements, two pen lifts; cited to KanoAI's hand-made Gujarati barakhadi templates. Finger-trace \textbf{\gu{ૌ}} once.

\subsection*{Writing — copy with the model}
Copy \textbf{\gu{સૌ}} once: \textbf{\gu{સ}}, the bar, then the two flags. Compare only two things: are there two flags, and does the lower one come first? Repair one if needed.

\subsection*{Writing — recall after hiding}
Hide the model and write the \emph{au} sign once. Uncover \textbf{\gu{ૌ}}, compare, and stop after one repair.

\begin{tcolorbox}[breakable,colback=teal!4,colframe=teal!35!black,title={Before you move on}]
Read \textbf{\gu{સો} \gu{સૌ} \gu{સૌથી}} aloud. What turns \textbf{\gu{ો}} into \textbf{\gu{ૌ}}? (\textbf{A second flag above.})

Sources: \href{https://github.com/gajjartejas/KanoAI/tree/9d3e2949a3f265448e430e15329093ea3aa516b8/interpolate-svg/svgs/barakhadi}{KanoAI's Gujarati barakhadi stroke templates} (Tejas Gajjar, KanoAI, Gujarati barakhadi centre-line stroke templates (interpolate-svg/svgs/barakhadi), commit 9d3e294, 1\_k/9\_kau.svg, groups g1, g2 then g3 (strokes g1s0, g2s0, g3s0): the au sign drawn after the consonant (GitHub, 2026)) and the \href{https://www.unicode.org/charts/PDF/U0A80.pdf}{Unicode Gujarati chart}.
\end{tcolorbox}
\section[\texorpdfstring{… \gu{કરતાં} \gu{મોટું} (… kartā̃ moṭũ)}{… કરતાં મોટું (… kartā̃ moṭũ)}]{… \gu{કરતાં} \gu{મોટું} (… kartā̃ moṭũ) — bigger than …}

\label{lesson:GU-C282-kartan}



\begin{quote}
Say \textquotedblleft{}Mira is the youngest\textquotedblright{}. (\emph{Mīrā sauthī nānī che}.) Then \textquotedblleft{}big\textquotedblright{} and \textquotedblleft{}a city\textquotedblright{}. (\textbf{\gu{મોટું}}, \textbf{\gu{શહેર}}.)
\end{quote}

\begin{grammarlens}[title={Grammar Lens: \gu{કરતાં}, than}]
\begin{quote}
\textbf{\gu{શહેર} \gu{ગામ} \gu{કરતાં} \gu{મોટું} \gu{છે}.}
\end{quote}

\begin{quote}
\emph{śaher gām kartā̃ moṭũ che.} — \textquotedblleft{}The city is bigger than the village.\textquotedblright{}
\end{quote}

The thing being described, then the one it is measured against with \textbf{\gu{કરતાં}} (\emph{kartā̃}), \textquotedblleft{}than\textquotedblright{}, then the describing word, which follows the first: \textbf{\gu{મીરા} \gu{અરુણ} \gu{કરતાં} \gu{નાની} \gu{છે}} (\emph{Mīrā Aruṇ kartā̃ nānī che}), \textquotedblleft{}Mira is younger than Arun\textquotedblright{}.

A word with no ending of its own takes \textbf{\gu{વધારે}}, \textquotedblleft{}more\textquotedblright{}, or \textbf{\gu{ઓછું}}, \textquotedblleft{}less\textquotedblright{}: \textbf{\gu{ટ્રેન} \gu{બસ} \gu{કરતાં} \gu{વધારે} \gu{ઝડપી} \gu{છે}} (\emph{ṭren bas kartā̃ vadhāre jhaḍpī che}).
\end{grammarlens}

\subsection*{You'll want to know: the most of all}
\begin{quote}
\textbf{\gu{ઝડપી} \gu{ટ્રેન} \gu{સૌથી} \gu{મોંઘી} \gu{છે}.}
\end{quote}

\begin{quote}
\emph{jhaḍpī ṭren sauthī moṅghī che.} — \textquotedblleft{}The fast train is the most expensive.\textquotedblright{}
\end{quote}

Two things take \textbf{\gu{કરતાં}}; the top of a whole group takes \textbf{\gu{સૌથી}}:

\noindent
\begin{tabularx}{\linewidth}{@{}>{\raggedright\arraybackslash}X>{\raggedright\arraybackslash}X@{}}
\toprule
\textbf{} & \textbf{say} \\
\midrule
the bus is cheaper than the train & \textbf{\gu{બસ} \gu{ટ્રેન} \gu{કરતાં} \gu{સસ્તી} \gu{છે}} \\
the bus is the cheapest & \textbf{\gu{બસ} \gu{સૌથી} \gu{સસ્તી} \gu{છે}} \\
\bottomrule
\end{tabularx}

\subsection*{Your turn}
Say these aloud:

\begin{itemize}
\raggedright
  \item \textquotedblleft{}Mira is younger than Arun\textquotedblright{} — \emph{Mīrā Aruṇ kartā̃ nānī che}
  \item \textquotedblleft{}The train is faster than the bus\textquotedblright{} — \emph{ṭren bas kartā̃ vadhāre jhaḍpī che}
\end{itemize}

\begin{tcolorbox}[breakable,colback=teal!4,colframe=teal!35!black,title={Before you move on}]
\textquotedblleft{}Than\textquotedblright{}? (\textbf{\gu{કરતાં}}.) \textquotedblleft{}The city is bigger than the village\textquotedblright{}? (\textbf{\gu{શહેર} \gu{ગામ} \gu{કરતાં} \gu{મોટું} \gu{છે}}.) For the top of a whole group? (\textbf{\gu{સૌથી}}.)
\end{tcolorbox}
\section[\texorpdfstring{\gu{સુધી} (sudhī)}{સુધી (sudhī)}]{\gu{સુધી} (sudhī) — until; as far as}

\label{lesson:GU-C282-sudhi}



\begin{quote}
Before the new one: say \textquotedblleft{}the city is bigger than the village\textquotedblright{}. (\emph{śaher gām kartā̃ moṭũ che}.) Then \textquotedblleft{}evening\textquotedblright{} and \textquotedblleft{}a river\textquotedblright{}. (\textbf{\gu{સાંજ}}, \textbf{\gu{નદી}}.)
\end{quote}

\subsection*{You'll want to know: \gu{સુધી}}
\textbf{\gu{સુધી}} — \emph{sudhī} — \textquotedblleft{}until; as far as\textquotedblright{}. It comes after its word, as \textbf{\gu{કરતાં}} does: an end in time, or in space.

\begin{quote}
\textbf{\gu{હું} \gu{સાંજ} \gu{સુધી} \gu{ઘરે} \gu{છું}.}
\end{quote}

\begin{quote}
\emph{hũ sānj sudhī ghare chũ.} — \textquotedblleft{}I'm at home until evening.\textquotedblright{}
\end{quote}

\begin{quote}
\textbf{\gu{અમે} \gu{નદી} \gu{સુધી} \gu{ચાલ્યા}.}
\end{quote}

\begin{quote}
\emph{ame nadī sudhī cālyā.} — \textquotedblleft{}We walked as far as the river.\textquotedblright{}
\end{quote}

\begin{cousinweb}
\textbf{\gu{સુધી}} (\emph{sudhī}) is a small word with no older form in any dictionary we use, so its history is not settled.
\end{cousinweb}

\subsection*{Your turn}
Say these aloud:

\begin{itemize}
\raggedright
  \item \emph{sudhī}, then \emph{sānj sudhī}
  \item \textquotedblleft{}I'm at home until evening\textquotedblright{} — \emph{hũ sānj sudhī ghare chũ}
  \item \textquotedblleft{}We walked as far as the river\textquotedblright{} — \emph{ame nadī sudhī cālyā}
\end{itemize}

\begin{tcolorbox}[breakable,colback=teal!4,colframe=teal!35!black,title={Before you move on}]
What does \textbf{\gu{સુધી}} mean? (\textquotedblleft{}Until; as far as\textquotedblright{}.) Before its word, or after? (After: \textbf{\gu{સાંજ} \gu{સુધી}}.)
\end{tcolorbox}
\section[\texorpdfstring{\gu{બે} \gu{વરસથી} (be varasthī)}{બે વરસથી (be varasthī)}]{\gu{બે} \gu{વરસથી} (be varasthī) — for how long, since when}

\label{lesson:GU-C282-be-varasthi}



\begin{quote}
Say \textquotedblleft{}until evening\textquotedblright{}. (\emph{sānj sudhī}.) Then \textquotedblleft{}a year\textquotedblright{} and \textquotedblleft{}morning\textquotedblright{}. (\textbf{\gu{વરસ}}, \textbf{\gu{સવાર}}.)
\end{quote}

\begin{grammarlens}[title={Grammar Lens: -\gu{થી} and the present}]
\begin{quote}
\textbf{\gu{હું} \gu{બે} \gu{વરસથી} \gu{અહીં} \gu{રહું} \gu{છું}.}
\end{quote}

\begin{quote}
\emph{hũ be varasthī ahī̃ rahũ chũ.} — \textquotedblleft{}I have lived here for two years.\textquotedblright{}
\end{quote}

\textbf{-\gu{થી}}, \textquotedblleft{}from\textquotedblright{}, on a length of time or a starting point says how long something has gone on, and the verb stays in the present: it is still going on. English says \textquotedblleft{}have lived\textquotedblright{}; Gujarati says \textquotedblleft{}I live, from two years\textquotedblright{}. \textbf{\gu{સવારથી} \gu{વરસાદ} \gu{પડે} \gu{છે}} (\emph{savārthī varsād paḍe che}): \textquotedblleft{}it has been raining since morning\textquotedblright{}. To ask: \textbf{\gu{તમે} \gu{કેટલા} \gu{વખતથી} \gu{અહીં} \gu{રહો} \gu{છો}?} (\emph{tame keṭlā vakhatthī ahī̃ raho cho?}), \textquotedblleft{}how long have you lived here?\textquotedblright{}
\end{grammarlens}

\subsection*{You'll want to know: from … until, and all day}
\begin{quote}
\textbf{\gu{દુકાન} \gu{સવારથી} \gu{સાંજ} \gu{સુધી} \gu{ખુલ્લી} \gu{છે}.}
\end{quote}

\begin{quote}
\emph{dukān savārthī sānj sudhī khullī che.} — \textquotedblleft{}The shop is open from morning until evening.\textquotedblright{}
\end{quote}

\textbf{-\gu{થી}} marks the start and \textbf{\gu{સુધી}} the end. A span with no start or end named takes \textbf{\gu{આખું}}, \textquotedblleft{}whole\textquotedblright{}, with no ending: \textbf{\gu{મેં} \gu{આખો} \gu{દિવસ} \gu{કામ} \gu{કર્યું}} (\emph{mẽ ākho divas kām karyũ}), \textquotedblleft{}I worked all day\textquotedblright{}.

\subsection*{Your turn}
Say these aloud:

\begin{itemize}
\raggedright
  \item \textquotedblleft{}How long have you lived here?\textquotedblright{} — \emph{tame keṭlā vakhatthī ahī̃ raho cho?}
  \item \textquotedblleft{}I worked all day\textquotedblright{} — \emph{mẽ ākho divas kām karyũ}
\end{itemize}

\begin{tcolorbox}[breakable,colback=teal!4,colframe=teal!35!black,title={Before you move on}]
\textquotedblleft{}For two years\textquotedblright{}? (\textbf{\gu{બે} \gu{વરસથી}}.) Past or present after it? (Present: \textbf{\gu{રહું} \gu{છું}}.) \textquotedblleft{}From morning until evening\textquotedblright{}? (\textbf{\gu{સવારથી} \gu{સાંજ} \gu{સુધી}}.)
\end{tcolorbox}
\section[\texorpdfstring{\gu{ઋતુ} (ṛtu)}{ઋતુ (ṛtu)}]{\gu{ઋતુ} (ṛtu) — a season}

\label{lesson:GU-C282-rtu}



\begin{quote}
Before the new one: say \textquotedblleft{}summer\textquotedblright{}, \textquotedblleft{}winter\textquotedblright{} and \textquotedblleft{}the monsoon\textquotedblright{}. (\textbf{\gu{ઉનાળો}}, \textbf{\gu{શિયાળો}}, \textbf{\gu{ચોમાસું}}.) Then \textbf{\gu{મોસમ}}: \textquotedblleft{}a season; the weather\textquotedblright{}.
\end{quote}

\subsection*{You'll want to know: \gu{ઋતુ}}
\textbf{\gu{ઋતુ}} — \emph{ṛtu} — \textquotedblleft{}a season\textquotedblright{}. It is feminine. \textbf{\gu{મોસમ}} is also the weather; \textbf{\gu{ઋતુ}} is one part of the year.

\begin{quote}
\textbf{\gu{અહીં} \gu{ત્રણ} \gu{ઋતુ} \gu{છે}: \gu{ઉનાળો}, \gu{ચોમાસું} \gu{અને} \gu{શિયાળો}.}
\end{quote}

\begin{quote}
\emph{ahī̃ traṇ ṛtu che: unāḷo, comāsũ ane śiyāḷo.} — \textquotedblleft{}There are three seasons here: summer, the monsoon and winter.\textquotedblright{}
\end{quote}

\begin{quote}
\textbf{\gu{તમને} \gu{કઈ} \gu{ઋતુ} \gu{સૌથી} \gu{વધારે} \gu{ગમે} \gu{છે}?}
\end{quote}

\begin{quote}
\emph{tamne kaī ṛtu sauthī vadhāre game che?} — \textquotedblleft{}Which season do you like most?\textquotedblright{}
\end{quote}

\textbf{\gu{કઈ}} is \textbf{\gu{કયું}}, \textquotedblleft{}which\textquotedblright{}, following \textbf{\gu{ઋતુ}}.

\begin{cousinweb}
\textbf{\gu{ઋતુ}} (\emph{ṛtu}) is a tatsama: Sanskrit \emph{ṛtu}, \textquotedblleft{}a fixed time; a season\textquotedblright{} (Monier-Williams), taken into Gujarati as it was, with the vowel letter \textbf{\gu{ઋ}} that the next lesson writes.
\end{cousinweb}

\subsection*{Your turn}
Say these aloud:

\begin{itemize}
\raggedright
  \item \emph{ṛtu}, then \emph{traṇ ṛtu}
  \item \textquotedblleft{}Which season do you like most?\textquotedblright{} — \emph{tamne kaī ṛtu sauthī vadhāre game che?}
  \item \textquotedblleft{}I like winter most of all\textquotedblright{} — \emph{mane śiyāḷo sauthī vadhāre game che}
\end{itemize}

\begin{tcolorbox}[breakable,colback=teal!4,colframe=teal!35!black,title={Before you move on}]
What does \textbf{\gu{ઋતુ}} mean? (\textquotedblleft{}A season\textquotedblright{}.) \textquotedblleft{}Which season?\textquotedblright{} (\textbf{\gu{કઈ} \gu{ઋતુ}?}: \textbf{\gu{ઋતુ}} is feminine.)
\end{tcolorbox}
\section[\texorpdfstring{\gu{ઋ} (ṛ)}{ઋ (ṛ)}]{\gu{ઋ} — the standing vowel ṛ}

\label{lesson:GU-W53-vocalic-r-letter}



\begin{quote}
\begin{itemize}
\raggedright
  \item \emph{Recall:} say \emph{a season} — \textbf{R1}, one lesson back
  \item \emph{Recall:} say \textbf{\gu{હૃદય}}, \textquotedblleft{}the heart\textquotedblright{}, and find the sign \textbf{\gu{ૃ}} under its \textbf{\gu{હ}}
\end{itemize}

\textbf{\gu{ઋતુ}} opens on the same vowel \emph{ṛ}, with no consonant to carry it.
\end{quote}

\subsection*{You'll want to know: \gu{ૃ}, standing alone}
\noindent
\begin{tabularx}{\linewidth}{@{}>{\raggedright\arraybackslash}X>{\raggedright\arraybackslash}X>{\raggedright\arraybackslash}X@{}}
\toprule
\textbf{} & \textbf{after a consonant} & \textbf{standing alone} \\
\midrule
\emph{ṛ} & \textbf{\gu{ૃ}}, as in \textbf{\gu{હૃદય}} & ? \\
\bottomrule
\end{tabularx}

The missing letter is \textbf{\gu{ઋ}}. Its print has three parts: a bent left body, a separate central stem with a foot, and a right loop flowing into a descending tail.

\subsection*{Script — meet one form}
\begin{quote}
\textbf{\gu{ઋ}} — \emph{ṛ}
\end{quote}

Keep \textbf{\gu{ઋ}} visible. Follow the numbered strip:

\begin{enumerate}
\raggedright
  \item sweep right along the upper left body and turn back down-left
  \item lift and descend the central stem into its foot
  \item lift again, circle the right loop, and finish through the descending tail
\end{enumerate}

Three strokes, two pen lifts: a verified three-stroke teaching-app variant fitted to the bundled printed outline (t30apps.com, ‘Gujarati Alphabet Writing Practice’, version 1.0, \gu{ઋ} animation, first through third SVG paths).

\subsection*{Writing — notice and trace}
\hlblockfigure{\detokenize{figures/GU-W53-vocalic-r-letter-filmstrip.pdf}}{How \gu{ઋ} is written, stroke by stroke}

Finger-trace \textbf{\gu{ઋ}} once, then air-trace it once. Notice the two lifts: before the stem, and before the loop.

\subsection*{Writing — copy with the model}
Copy visible \textbf{\gu{ઋ}} once. Compare only two things: are there three parts, and does the loop come last? Repair one landmark if needed.

\subsection*{Writing — recall after hiding}
Hide the model and write the standing vowel \emph{ṛ} once. Uncover \textbf{\gu{ઋ}}, compare, and stop after one repair.

\begin{tcolorbox}[breakable,colback=teal!4,colframe=teal!35!black,title={Before you move on}]
Read \textbf{\gu{ૃ} \gu{ઋ} \gu{ઋતુ}} aloud. Which part of \textbf{\gu{ઋ}} is written last? (\textbf{The right loop and its tail.})

Sources: \href{https://www.t30apps.com/gujarati-alphabet-writing/}{Gujarati Alphabet Writing Practice} and the \href{https://www.unicode.org/charts/PDF/U0A80.pdf}{Unicode Gujarati chart}.
\end{tcolorbox}
\section[(review)]{Review — comparing, and how long}

\label{lesson:GU-R282-comparing-and-time}



\begin{quote}
Say \textquotedblleft{}I'm afraid\textquotedblright{}. (\emph{mane bīk lāge che}.) Then \textquotedblleft{}I want to become a doctor\textquotedblright{}. (\emph{māre dāktar banvũ che}.) And \textquotedblleft{}my dream is to become a doctor\textquotedblright{}. (\emph{mārũ sapnũ dāktar banvānũ che}.)

\begin{itemize}
\raggedright
  \item \emph{Write it:} \textbf{\gu{ૌ}} in \textbf{\gu{સૌથી}}, then \textbf{\gu{ઋ}} in \textbf{\gu{ઋતુ}}, from memory
\end{itemize}
\end{quote}

\subsection*{Your turn}
Say these aloud:

\begin{itemize}
\raggedright
  \item The city is bigger than the village. Mira is younger than Arun.
  \item The bus is the cheapest. The fast train is the most expensive.
  \item I have lived here for two years. It has been raining since morning.
  \item The shop is open from morning until evening. I worked all day.
  \item Which season do you like most? I like the monsoon most of all.
\end{itemize}

\begin{tcolorbox}[breakable,colback=teal!4,colframe=teal!35!black,title={Before you move on}]
\textquotedblleft{}Than\textquotedblright{}? (\textbf{\gu{કરતાં}}.) \textquotedblleft{}Most of all\textquotedblright{}? (\textbf{\gu{સૌથી} \gu{વધારે}}.) \textquotedblleft{}For two years\textquotedblright{}: which tense follows? (The present.) Which sign is on \textbf{\gu{સ}} in \textbf{\gu{સૌથી}}? (\textbf{\gu{ૌ}}.)
\end{tcolorbox}
